\documentclass[11pt]{article}
\usepackage[margin=1in]{geometry}
\usepackage{amsmath,amssymb,amsthm}
\usepackage[hidelinks]{hyperref}
\newtheorem{theorem}{Theorem}
\title{Sharp zero-free gaps for\\half-integer steps in three-frequency progressions}
\author{}
\date{28 November 2026}
\begin{document}
\maketitle
\vspace{-2em}
\begin{abstract}
We determine the longest zero-free interval, as a supremum, for
every real trigonometric spectrum \(\{\pm 2n,\pm(d+2)n,\pm(2d+2)n\}\)
with \(n\ge1\) and odd \(d\ge5\). There are two analytic branches,
one for each residue of \(d\bmod4\). For every spectrum the critical
closed interval has a unique nonnegative polynomial up to positive
scale and a unique annihilating probability measure, necessarily
on six points. A replayable rational certificate proves the two
branches on their full parameter ranges.
\end{abstract}

\section{Exact answer}
Kozma and Oravecz ask for the sharp zero-free length of real
trigonometric polynomials with frequencies in a symmetric arithmetic
progression of step at least twice its start \cite[Question~1]{KO}.
The positive-measure formulation in ProveIt \cite{P} gives dual
certificates for this extremum. This note treats three frequencies when
the step-to-start ratio is a nonintegral half-integer. The integer-ratio
case is covered in the companion note. Neither uses RH or GRH.

Set \(r=d/2\), \(\epsilon=1/r\), and
\(\sigma=-1\) if \(d\equiv1\pmod4\), \(+1\) if \(d\equiv3\pmod4\).
For \(\epsilon\ge 0\) put \(h_\epsilon(s)=\sin(\epsilon s)/\epsilon\),
with \(h_0(s)=s\), and \(L_\epsilon=\pi/(2\epsilon)\) with \(L_0=+\infty\).
The other two functions are
\[
 C_{\sigma,\epsilon}(s)=\cos((1+\epsilon)s)+\sigma\sin((1+\epsilon)s),
 \qquad D_\epsilon(s)=\cos((2+\epsilon)s).
\]
For each sign let \(A,B,u,v,x=x_\sigma(\epsilon)\) be the certified
solution to the equations
\begin{gather*}
 F=h_\epsilon+A C_{\sigma,\epsilon}+B D_\epsilon,\\
 F(u)=F(v)=F'(u)=F'(v)=F(x)=0,\qquad -2<x<0<u<v<L_\epsilon.
\end{gather*}
Here and below derivatives of \(F\) with respect to \(s\) are marked
with primes. The inequalities in Section~\ref{sec:interval} specify
the branch uniquely. Its enclosure can be obtained at any parameter
by locating the appropriate row in the bundled certificate.

\begin{theorem}\label{thm:half}
Let \(M\) be the supremal zero-free length on \(\mathbb R/\mathbb Z\)
for the spectrum \(\{\pm 2n,\pm(d+2)n,\pm(2d+2)n\}\), where \(d\ge5\)
is odd and \(n\) is a positive integer. Then
\[
 M=\frac1{4n}-\frac{x_\sigma(2/d)}{\pi dn}.
\]
The supremum is not attained even on an open interval. On the fixed
critical closed interval there is a unique ray of nonnegative
polynomials and a unique annihilating probability measure,
with six atoms. If \(N=2n\), then for each residue class as \(r\to\infty\),
\[
 NM=\frac12+\frac{c_{1,\sigma}}{r}+
               \frac{c_{2,\sigma}}{r^2}+O(r^{-3}),\qquad
 c_{1,\sigma}=-\frac{x_\sigma(0)}{\pi},\quad
 c_{2,\sigma}=-\frac{x_\sigma'(0)}{\pi}.
\]
\end{theorem}

\section{Two compact interval certificates}\label{sec:interval}
Use the parameter ranges \(0\le\epsilon\le 2/5\) for \(\sigma=-1\)
and \(0\le\epsilon\le 2/7\) for \(\sigma=+1\).
The replayable checks establish, uniformly on these ranges:
\begin{enumerate}
\item The five equations have a solution in the indicated order,
with nonsingular Jacobian.
On \([x,L_\epsilon]\), \(F\ge0\) with zeros precisely at \(x,u,v\).
The derivatives \(F'(x),F''(u),F''(v)\) are strictly positive.
\item With \(w(s)=(C_{\sigma,\epsilon}(s),D_\epsilon(s))\), the
three cyclic determinants of \(w(x),w(u),w(v)\) are nonzero with the
same sign.
\item The three rows
\[
 \bigl(\cos(\epsilon s),\
   \sin((1+\epsilon)s)-\sigma\cos((1+\epsilon)s),\
   \sin((2+\epsilon)s)\bigr),\quad s=x,u,v,
\]
form an invertible matrix.
\end{enumerate}

\paragraph{Root boxes.}
Write \(z=(A,B,u,v,x)\) and
\(g=(F(u),F(v),F'(u),F'(v),F(x))\).
For each parameter interval \(E\), the data contain a center \(z_0\),
a radius \(a>0\), and a numerical inverse proposal converted to a
rational matrix \(R\). Put \(Z=z_0+[-a,a]^5\).
Enclose \(Rg(E,z_0)\) and \(I-R D_z g(E,Z)\).
With \(q_i\) the sum of the absolute-value upper bounds in row \(i\)
of the second enclosure, check
\[
 q_i<1,\qquad
 \sup_{\epsilon\in E}|(Rg(\epsilon,z_0))_i|+a q_i<a.
\]
Then \(z\mapsto z-Rg(\epsilon,z)\) is a contraction taking \(Z\)
strictly inside itself for each \(\epsilon\in E\). The norm bound
implies \(R D_zg\) is invertible, so \(R\) and \(D_zg\) are invertible
and the fixed point solves \(g=0\). Completeness of the full parameter
cover is checked by equality of adjacent rational endpoints.

\paragraph{Positivity.}
For each root box the order \(-2<x<0<u<v<L_\epsilon\) is checked,
using \(3.14<\pi\) for the last inequality.
Enlarge the three root-coordinate intervals by \(1/9\) rounded down
to the arithmetic grid. On the enlarged \(x\) interval check \(F'>0\);
on those for \(u,v\) check \(F''>0\). This gives the signs near the
zeros. The intervening intervals and the portion after \(v\) up
to the compact cutoff are checked by subdivision and direct
positive lower bounds on \(F\). The cutoff is 9 for a parameter
interval with lower endpoint zero, otherwise
\(\min(9,11/(7\epsilon_-))\), rounded up to the grid.
Since \(\pi<22/7\), this covers up to \(\min(9,L_\epsilon)\).
The checked bound \(3|A|+2|B|<126/11\) and
\(|C_{\sigma,\epsilon}|\le\sqrt2<3/2\) imply that for the tail
\(9\le s\le L_\epsilon\),
\[
 h_\epsilon(s)\ge 2s/\pi\ge18/\pi>63/11>
                  \sqrt2\,|A|+|B| .
\]
The cyclic determinants and sine determinant are evaluated on the
root boxes with the same arithmetic.

\paragraph{Outward rounding.}
Every sign is checked on a \(10^{-23}\) fixed grid using integer
floor and ceiling. For a sine or cosine interval, evaluate the midpoint
divided by eight by Taylor polynomials of degrees 43 or 42 respectively,
and use three double-angle steps. The initial reduced arguments are
checked to lie in \([-3,3]\); \(3^{44}/44!<10^{-23}\) bounds the Taylor
remainder. The same bound applies to the integrated cosine series for
\(\sin y/y=\int_0^1\cos(ty)\,dt\) on this reduced domain. Input radii
are propagated by the Lipschitz bounds \(1,1,1/2\), respectively.
The final covers have 354 and 780 cells and use 70,624 and 77,808
compact sign evaluations. The replay uses no floating-point
proposals; a separate optional mode regenerates the rational data.

\section{Sharpness and rigidity}
For a fixed parameter, take the three cyclic determinants in
item 2 and multiply by a common sign to make positive weights
\(\lambda_x,\lambda_u,\lambda_v\).
The \(2\times 3\) determinant identity gives
\(\sum_s\lambda_s w(s)=0\); because \(F(x)=F(u)=F(v)=0\),
the same weights cancel \(h_\epsilon\).

This also proves uniqueness of the branch specified above.
If two such endpoints satisfied \(x_1<x_2<0\), the first candidate
would be nonnegative on the second contact set and strictly positive
at \(x_2\), contradicting the second weights.
At equal endpoints the weights force equality of the two contact
sets. The equations at \(u,v\) then fix \(A,B\) by their nonzero
minor. The analytic implicit function theorem applies on every
cell; uniqueness at common endpoints glues the analytic branches.
In particular each has a germ at \(\epsilon=0\).

Put \(\theta=2\pi(2n)t\) and \(s=r(\pi/2-\theta)\). By expanding
the phases at \(\theta=\pi/2\), the three cosine modes at
\(1,1+r,1+2r\) in \(\theta\) become nonzero scalar multiples of
\(h_\epsilon,C_{\sigma,\epsilon},D_\epsilon\), respectively.
Indeed \(\tan((1+r)\pi/2)=\sigma\), and \(1+2r\) is even.
The critical half-angle is \(\theta_*=\pi/2-x/r\).
Place the weights \(\lambda_x,\lambda_u,\lambda_v\) in equal pairs
at angles \(\pm(\pi/2-s/r)\), \(s=x,u,v\), and normalize.
The cosine moments vanish by construction and the sine moments
by reflection.

If a polynomial were zero-free on an interval of angular length
\(2\theta_*\), translate its center to zero and adjust the sign.
It is strictly positive in the interior, and nonnegative at the
endpoints by continuity. Its integral against the paired measure
is strictly positive (four atoms are interior), a contradiction.
A larger interval contains one of the critical length and is
likewise impossible. For any \(x<a<0\), the combination
\(F+\delta h_\epsilon\) is positive for sufficiently small
\(\delta>0\) on \([a,L_\epsilon]\). On \([a,0]\) this follows
from the positive minimum of \(F\); above zero the added term
is positive, including at both double zeros. The resulting
interval lengths approach the critical value. Conversion from
\(\theta\) to \(t\) gives \(M=\theta_*/(2\pi n)\).

If any polynomial is nonnegative throughout the critical interval,
integration shows it must vanish at all six points. The odd part
vanishes by the determinant in item 3: after expansion at \(\pi/2\),
its three columns are nonzero scalar multiples of the sine-mode
evaluation columns. The even equations at \(u,v\) have rank two,
so the even part is a multiple of the certified combination.
The multiple must be nonnegative. Conversely, integrating that
combination against any annihilating probability measure forces
its support onto the six points. The odd determinant makes the
pair masses equal; the cosine equations have rank two on the three
pair totals. Normalization determines their unique positive values,
so no atom can be deleted.

\section{Asymptotics}
The analytic germ gives the asymptotic expansion in
Theorem~\ref{thm:half}. Its derivatives can be computed by a
five-by-five linear solve with the same Jacobian as above:
\(D_zg\,z'(0)=-\partial_\epsilon g\) at zero.
For example, holding coefficients fixed, the first derivative of
\(F\) in \(\epsilon\) at zero is
\[
 A\,s(-\sin s+\sigma\cos s)-B\,s\sin(2s).
\]
Approximate constants are \(c_{1,-}=0.6064836353\),
\(c_{2,-}=-0.3704115178\), \(c_{1,+}=0.2224839370\), and
\(c_{2,+}=-0.1254621878\). The exact definitions are the germ
expressions; these decimal approximations are not used by the proof.
The distinct \(1/r\) coefficients show that a single expansion
without the residue class cannot describe both branches.

\clearpage
\section{Independent first case and numerical comparison}
For \(d=5\) there is also a separate algebraic certificate.
Use \(T_j\) for Chebyshev polynomials of the first kind and set
\(P=T_2+A_0T_7+B_0T_{12}\). There is a unique triple with
\(P(u_0)=P(v_0)=P'(u_0)=P'(v_0)=0\) after eliminating \(A_0\),
in the specified boxes
\begin{align*}
 B_0&\in[0.46688711936350,\ 0.46688711936362],\\
 u_0&\in[0.94228289217292,\ 0.94228289217299],\\
 v_0&\in[0.72954822734434,\ 0.72954822734442].
\end{align*}
The checker obtains the two stationary points of
\(R=(T_2+B_0T_{12})/T_7\) using opposite endpoint signs
and monotonicity of the numerator of \(R'\) in each box.
Then \(R(u_0,B_0)-R(v_0,B_0)\) has derivative
\((T_{12}/T_7)(u_0)-(T_{12}/T_7)(v_0)>0\) and opposite
signs at the ends of the \(B_0\) interval.
Set \(A_0=-R(u_0,B_0)\). The remaining endpoint \(\xi\)
is the unique simple root in \((0.46362209927,0.46362209930)\).
After factoring \((y-\xi)(y-u_0)^2(y-v_0)^2\) out of \(P(y)\),
the degree-seven quotient has 256 positive Bernstein bounds
across 32 equal subintervals of \([0.46362209927,1]\).
The signs of \(T_7\) at \(\xi,u_0,v_0\) are \(+,-,+\),
with \(T_{12}/T_7\) at \(u_0\) strictly between the other two
ratios. This gives the same positive dual argument independently,
and \(T_2\) is positive at the two double roots, so it permits
the perturbation for the lower bound. Thus
\(M=\arccos(\xi)/(\pi n)\simeq0.346549538462/n\).
Here \(y=\cos(2\pi nt)\), unlike the angular rescaling in the main
certificate. The exact sine determinant is also checked separately.

For comparison, here are approximate scaled lengths from the two
branches:
\[
\begin{array}{c|rrrr}
d & 5 & 7 & 9 & 11\\ \hline
n M & 0.346549538462 & 0.276881103951 &
      0.309113910011 & 0.268194148020
\end{array}
\]
The result addresses the wide three-frequency part of Question~1
for half-integer ratios. It does not classify further denominators
or progressions with more than three frequencies.

\begin{thebibliography}{9}
\bibitem{KO} G. Kozma and F. Oravecz,
\emph{On the gaps between zeros of trigonometric polynomials},
Real Anal. Exchange 28 (2002/03), no.~2, 447--454.
\href{https://arxiv.org/abs/math/0601021}{arXiv:math/0601021}.
\bibitem{P} ProveIt,
\href{https://github.com/VladimirReshetnikov/ProveIt/blob/3c4d27a69c98d09319a51e41a19fca05bb80a3e2/Analysis/GeneralizedRiemannHypothesis/Article/sections/fil.tex}{GRH article, spectral-filter chapter},
revision \texttt{3c4d27a69c98}.
\end{thebibliography}
\end{document}
